\section{Unique Paths}
\label{sec:dp-unique-paths}

\problemstatement{A robot starts at the top-left corner of an
$m \times n$ grid and wants to reach the bottom-right corner, moving only
\textbf{right} or \textbf{down} at each step. Count the number of unique
paths.}

\begin{patternbox}
\emph{``Number of paths through a grid, moving only right or down''} is
the cleanest possible grid-DP recurrence: a path arrives at cell
$(i,j)$ from exactly one of two neighbors -- directly above, $(i-1,j)$,
or directly to the left, $(i,j-1)$ -- and since every path to $(i,j)$ is
either one of these two cases, the count of paths to $(i,j)$ is their
\textbf{sum}, exactly mirroring the DP 1D chapter's Climbing Stairs
``sum over a fixed set of predecessors,'' just with the predecessors now
arranged on two grid axes instead of one line.
\end{patternbox}

\subsection*{Deriving the recurrence}

Let $\textit{ways}(i,j)$ be the number of unique paths from $(0,0)$ to
$(i,j)$. Since the only moves are right and down, the last move into
$(i,j)$ came from $(i-1,j)$ (a down-move) or $(i,j-1)$ (a right-move):
\[
  \textit{ways}(i,j) = \textit{ways}(i-1,j) + \textit{ways}(i,j-1).
\]
Base case: $\textit{ways}(0,0) = 1$ (the trivial path of staying put);
any $\textit{ways}(i,j)$ with $i<0$ or $j<0$ is treated as $0$
(non-existent, out-of-grid predecessors contribute nothing).

\subsection*{Approach 1: Recursion}

\begin{lstlisting}
#include <bits/stdc++.h>
using namespace std;

int uniquePathsRecursive(int i, int j) {
    if (i == 0 && j == 0) return 1;
    if (i < 0 || j < 0) return 0;

    int fromAbove = uniquePathsRecursive(i - 1, j);
    int fromLeft  = uniquePathsRecursive(i, j - 1);
    return fromAbove + fromLeft;
}

int main() {
    int m = 3, n = 7;
    cout << uniquePathsRecursive(m - 1, n - 1) << endl; // 28
    return 0;
}
\end{lstlisting}

\begin{complexitybox}[Approach 1 Complexity]
\textbf{Time.} $O(2^{m+n})$ -- each call branches into two, with
recursion depth up to $m+n$.
\textbf{Space.} $O(m+n)$ for the recursion stack.
\end{complexitybox}

\subsection*{Approach 2: Memoization}

\begin{lstlisting}
#include <bits/stdc++.h>
using namespace std;

int uniquePathsMemo(int i, int j, vector<vector<int>>& memo) {
    if (i == 0 && j == 0) return 1;
    if (i < 0 || j < 0) return 0;
    if (memo[i][j] != -1) return memo[i][j];

    int fromAbove = uniquePathsMemo(i - 1, j, memo);
    int fromLeft  = uniquePathsMemo(i, j - 1, memo);
    return memo[i][j] = fromAbove + fromLeft;
}

int main() {
    int m = 3, n = 7;
    vector<vector<int>> memo(m, vector<int>(n, -1));
    cout << uniquePathsMemo(m - 1, n - 1, memo) << endl; // 28
    return 0;
}
\end{lstlisting}

\begin{complexitybox}[Approach 2 Complexity]
\textbf{Time.} $O(m \cdot n)$ -- one $O(1)$ computation per distinct
$(i,j)$ pair.
\textbf{Space.} $O(m \cdot n)$ memo $+$ $O(m+n)$ recursion stack.
\end{complexitybox}

\subsection*{Approach 3: Tabulation}

\begin{lstlisting}
#include <bits/stdc++.h>
using namespace std;

int uniquePathsTabulation(int m, int n) {
    vector<vector<int>> dp(m, vector<int>(n, 0));

    for (int i = 0; i < m; i++) {
        for (int j = 0; j < n; j++) {
            if (i == 0 && j == 0) {
                dp[i][j] = 1;
                continue;
            }
            int fromAbove = (i > 0) ? dp[i-1][j] : 0;
            int fromLeft  = (j > 0) ? dp[i][j-1] : 0;
            dp[i][j] = fromAbove + fromLeft;
        }
    }
    return dp[m - 1][n - 1];
}

int main() {
    cout << uniquePathsTabulation(3, 7) << endl; // 28
    return 0;
}
\end{lstlisting}

\subsection*{Dry run of the tabulation table}

\dryrun{Take $m=3$, $n=3$.}

\begin{center}
\begin{tabular}{c|c|c|c}
 & $j=0$ & $j=1$ & $j=2$ \\ \hline
$i=0$ & 1 & 1 & 1 \\ \hline
$i=1$ & 1 & 2 & 3 \\ \hline
$i=2$ & 1 & 3 & 6
\end{tabular}
\end{center}

The first row and first column are all $1$s (only one way to reach any
cell along the top edge or left edge: a straight line of moves). Every
other cell is the sum of its top and left neighbors, e.g.\ $dp[2][2] =
dp[1][2]+dp[2][1] = 3+3=6$. The answer for a $3\times3$ grid is $6$.

\begin{complexitybox}[Approach 3 Complexity]
\textbf{Time.} $O(m \cdot n)$.
\textbf{Space.} $O(m \cdot n)$ table, no recursion stack.
\end{complexitybox}

\subsection*{Approach 4: Space Optimization}

\begin{ideabox}[Only the Previous Row Is Ever Needed]
Computing row $i$ only reads from row $i-1$ (directly above) and from
\emph{within} row $i$ itself (directly left, already computed earlier in
the same row's left-to-right scan). No row before $i-1$ is ever needed
again once row $i$ is being filled, so a single 1D array, reused and
updated in place across rows, suffices.
\end{ideabox}

\begin{lstlisting}
#include <bits/stdc++.h>
using namespace std;

int uniquePathsSpaceOptimized(int m, int n) {
    vector<int> prevRow(n, 0);

    for (int i = 0; i < m; i++) {
        vector<int> currentRow(n, 0);
        for (int j = 0; j < n; j++) {
            if (i == 0 && j == 0) {
                currentRow[j] = 1;
                continue;
            }
            int fromAbove = prevRow[j];
            int fromLeft  = (j > 0) ? currentRow[j-1] : 0;
            currentRow[j] = fromAbove + fromLeft;
        }
        prevRow = currentRow;
    }
    return prevRow[n - 1];
}

int main() {
    cout << uniquePathsSpaceOptimized(3, 7) << endl; // 28
    return 0;
}
\end{lstlisting}

\begin{complexitybox}[Approach 4 Complexity]
\textbf{Time.} $O(m \cdot n)$.
\textbf{Space.} $O(n)$ for the single rolling row, instead of
$O(m \cdot n)$ for the full table.
\end{complexitybox}

\begin{takeawaybox}
Grid DP's space optimization mirrors 1D DP's exactly, just one dimension
up: instead of two rolling \emph{scalars} (Fibonacci, House Robber), we
keep two rolling \emph{rows} (or, with more care, overwrite a single row
in place, since the left-neighbor read within a row happens before that
position is overwritten in a left-to-right scan). Whenever a grid
recurrence only reads from the immediately preceding row and the current
row, this reduction from $O(mn)$ to $O(n)$ space is always available.
\end{takeawaybox}
